\documentclass[conference]{IEEEtran}
\usepackage{times}

\usepackage[numbers]{natbib}
\usepackage{multicol}
\usepackage{graphicx}
\usepackage{booktabs}
\usepackage{xcolor}

\usepackage[bookmarks=true]{hyperref}

\pdfinfo{
   /Author (Anonymous)
   /Title  (AutoDex Rebuttal)
   /CreationDate (D:20260325120000)
   /Subject (Robotics)
   /Keywords (Dexterous;Grasping)
}

\definecolor{R1}{RGB}{0,128,0}
\definecolor{R2}{RGB}{0,0,200}
\definecolor{R3}{RGB}{200,0,0}
\definecolor{AC}{RGB}{128,0,128}

\newcommand{\rone}{{\color{R1}R1}}
\newcommand{\rtwo}{{\color{R2}R2}}
\newcommand{\rthree}{{\color{R3}R3}}
\newcommand{\acmark}{{\color{AC}AC}}

\newcommand{\rssparagraph}[1]{{\bf #1.}}

\begin{document}

\title{AutoDex [Paper-ID: 99]\vspace{-1cm}}
\maketitle
\thispagestyle{empty}

We thank the reviewers for their thoughtful feedback:
{\color{R1}R1} recognized our ``fully automated pipeline'' and ``compelling constrained scenarios'';
{\color{R2}R2} valued the ``clear quantitative argument for real-world filtering (33\%$\to$71\%)'' and ``coverage-preserving selection design'';
{\color{R3}R3} acknowledged the ``important problem'' and ``inclusion of failure cases.''
We address all concerns below.

\rssparagraph{[\rone, \rthree, \acmark] The necessity of real-world validation}
We do not position AutoDex as competing with sim DR --- rather, it addresses what DR fundamentally cannot. DR randomizes known parameters (friction, mass), but failure modes such as finger-pad deformation, cable backlash, and non-linear contact are not parameterized in current simulators.
To verify this, we tested [N] grasps under both DR filtering ([X] randomized conditions: friction $\times$[range], mass $\times$[range]) and real-robot execution.

\begin{table}[h]
\centering
\small
\caption{Sim-DR vs Real-World Outcomes ($N$=[X] grasps)}
\label{tab:dr_vs_real}
\begin{tabular}{l|cccc}
\toprule
\textbf{Object} & \textbf{DR Succ} & \textbf{Real Succ} & \textbf{False Pos} & \textbf{False Neg} \\ \midrule
Obj-1  & {[X]}\% & {[X]}\% & {[X]}\% & {[X]}\% \\
Obj-2  & {[X]}\% & {[X]}\% & {[X]}\% & {[X]}\% \\
Obj-3  & {[X]}\% & {[X]}\% & {[X]}\% & {[X]}\% \\
Obj-4  & {[X]}\% & {[X]}\% & {[X]}\% & {[X]}\% \\
Obj-5  & {[X]}\% & {[X]}\% & {[X]}\% & {[X]}\% \\
Obj-6  & {[X]}\% & {[X]}\% & {[X]}\% & {[X]}\% \\
Obj-7  & {[X]}\% & {[X]}\% & {[X]}\% & {[X]}\% \\
Obj-8  & {[X]}\% & {[X]}\% & {[X]}\% & {[X]}\% \\
Obj-9  & {[X]}\% & {[X]}\% & {[X]}\% & {[X]}\% \\
Obj-10 & {[X]}\% & {[X]}\% & {[X]}\% & {[X]}\% \\
\midrule
Aggregate & {[X]}\% & {[X]}\% & {[X]}\% & {[X]}\% \\
\bottomrule
\end{tabular}

{\small McNemar's $p < $ {[X]}, Cohen's $\kappa$ = {[X]}}
\end{table}

Of DR-robust grasps, [FP]\% failed on the real robot (false positives) --- passing every sim condition yet failing due to unmodeled dynamics. Conversely, [FN]\% were rejected by DR but succeeded in reality (false negatives) --- valid grasps that DR permanently discards.

\noindent[\rthree, \acmark] We use sim contact modeling not as ground truth, but to filter obvious failures --- grasps with no object contact or mesh penetration --- under deliberately optimistic settings (high friction, low gravity) to preserve borderline candidates rather than discard them. Real execution then filters out the false positives that this coarse check misses.

\rssparagraph{[\rtwo, \rthree, \acmark] Strengthened evidence for physical validity}
We expanded our evaluation to [15] objects (5 new objects spanning metal, thin, large categories) with [20] trials per cell.

\begin{table}[h]
\centering
\small
\caption{Expanded Physical Validation (Success Rate $\pm$ 95\% CI, $n$=[X] per cell)}
\label{tab:expanded_validation}
\begin{tabular}{ll|cc}
\toprule
\textbf{Cat.} & \textbf{Object} & \textbf{Sim-val.} & \textbf{Real-filt.} \\ \midrule
Tabletop & Obj-9  & {[X]}$\pm${[X]} & {[X]}$\pm${[X]} \\
         & Obj-3  & {[X]}$\pm${[X]} & {[X]}$\pm${[X]} \\
         & Obj-4  & {[X]}$\pm${[X]} & {[X]}$\pm${[X]} \\
         & Obj-5  & {[X]}$\pm${[X]} & {[X]}$\pm${[X]} \\
         & Obj-10 & {[X]}$\pm${[X]} & {[X]}$\pm${[X]} \\
\midrule
Complex  & Obj-1  & {[X]}$\pm${[X]} & {[X]}$\pm${[X]} \\
         & Obj-2  & {[X]}$\pm${[X]} & {[X]}$\pm${[X]} \\
         & Obj-3  & {[X]}$\pm${[X]} & {[X]}$\pm${[X]} \\
         & Obj-4  & {[X]}$\pm${[X]} & {[X]}$\pm${[X]} \\
         & Obj-7  & {[X]}$\pm${[X]} & {[X]}$\pm${[X]} \\
\midrule
New      & Obj-A  & {[X]}$\pm${[X]} & {[X]}$\pm${[X]} \\
         & Obj-B  & {[X]}$\pm${[X]} & {[X]}$\pm${[X]} \\
         & Obj-C  & {[X]}$\pm${[X]} & {[X]}$\pm${[X]} \\
         & Obj-D  & {[X]}$\pm${[X]} & {[X]}$\pm${[X]} \\
         & Obj-E  & {[X]}$\pm${[X]} & {[X]}$\pm${[X]} \\
\midrule
\multicolumn{2}{l|}{Aggregate} & {[X]}$\pm${[X]} & {[X]}$\pm${[X]} \\
\bottomrule
\end{tabular}

{\small Fisher's exact test (aggregate): $p < $ {[X]}}
\end{table}

The expanded evaluation confirms statistically significant improvement from physical filtering across diverse object categories (Fisher's exact, $p < $ [X]). Real-filtered grasps achieve [X]\% $\pm$ [X] vs [X]\% $\pm$ [X] for sim-validated, with consistent gains across tabletop, complex, and newly added object categories.

\rssparagraph{[\rthree, \acmark] Scalability of the multi-camera setup}
The pipeline is modular: capture, calibration, and evaluation are decoupled, so adapting to a different robot-hand or camera setup requires only modifying the relevant module. Calibration takes $\sim$40 board images across all cameras simultaneously ($\sim$10 min) plus automated hand-eye calibration ($\sim$10 min).
The 24-camera configuration is not a pipeline requirement but a research choice to study the effect of view density on grasp detection. Our pipeline operates with as few as 2 RGB cameras (or a single depth camera), though pose estimation accuracy degrades --- as shown in Table III, where a 4-cam setup increases ADD-S error from 0.31mm to 1.17mm and drops success from 62\% to 51\%. Dense multi-view does help --- during grasping, the hand occludes an average of [X]\% of the object --- but the camera count is a trade-off each lab makes based on their budget and goals.

\rssparagraph{[\rtwo, \acmark] Efficiency compared to teleoperation}
AutoDex achieves [X] grasps/hr vs [X] grasps/hr for teleop, requiring only [X] min of human time per grasp (object placement and reset). Teleop requires continuous operator presence --- wearing a body suit, real-time 22-DoF control --- with practical sessions limited to [X] hours before fatigue and Xsens charging. This makes 24/7 collection impractical. AutoDex runs unattended after object placement, enabling overnight collection --- and scales linearly with additional robots without additional operators. For dataset capture, infrastructure costs (robot, cameras) are shared regardless of collection method; teleop additionally requires Xsens MVN Awinda (\$[X]) with per-session setup overhead (suit calibration, sensor charging).

\rssparagraph{[\rtwo, \rthree] Grasp selection and scenario coverage}
\noindent[\rthree] Our procedural constraints systematically cover the combinatorial space of planar boundaries around an object: one side wall = Wall; one side + ceiling = Shelf; 2--3 sides with/without ceiling = Shelf variants; 4 sides open-top = Box. The only missing case --- fully enclosed --- makes grasping infeasible for any method.

\noindent[\rtwo] To confirm that sim solvability gains translate to real, we evaluated both sim and real solvability up to $N$=50, and extended sim evaluation to $N$=5000 to find the ceiling.

\begin{figure}[h]
\centering
% \includegraphics[width=\columnwidth]{figures/solvability_curve.pdf}
\fbox{\parbox{0.9\columnwidth}{\centering\vspace{1.5cm}[PLACEHOLDER: Solvability curve]\vspace{1.5cm}}}
\caption{Scenario solvability vs grasp budget $N$. Solid: sim, dashed: real ($N \leq 50$). Lines: Baseline sim/real, Ours (Greedy) sim/real, Ours (Random) sim. Ceiling converges at $\sim$[X]\%.}
\label{fig:solvability_curve}
\end{figure}

\noindent[\rtwo] Ours maintains higher real solvability than Baseline at $N$=10 and $N$=50, confirming that sim solvability gains are preserved in physical execution.
The sim curve converges at $\sim$[X]\% around $N$=[X]. Of the remaining unsolved scenarios, [X]\% are kinematically infeasible --- the Allegro hand cannot reach the object from any direction due to obstacle density --- and only [X]\% are coverage gaps addressable by generating more candidates.
As shown in Table I, greedy gain is largest at our actual validation budget ($N$=100: [X]pp, Table I Row 4 vs Row 3). At large $N$, random sampling achieves sufficient coverage through volume, so convergence is expected.

\rssparagraph{[\acmark, \rone, \rthree] Scope and contributions}
\noindent[\acmark, \rthree] We clarify AutoDex's contribution beyond engineering. While the real-world validation pipeline itself is an engineering effort, it enables something previously unavailable: physically-grounded grasp data free from human demonstration bias. Teleoperation datasets are inherently biased toward operator-preferred approach paths; AutoDex decouples grasp generation from execution, allowing any generation method (BODex, RL policies, human priors) to be validated on real hardware. This generator-agnostic design is a methodological contribution, not a single-use BODex extension.
Beyond automation, we address the core scalability bottleneck of real validation: the coverage-preserving selection algorithm maximizes the information gained per costly real execution, ensuring diverse constraint coverage within a limited budget. This is not simply ``running more grasps'' --- it formalizes \emph{which} grasps to validate as a combinatorial optimization problem.

\noindent[\rone, \acmark] The current dataset focuses on stable grasps --- an intentional scope decision. Stable grasping under environmental constraints remains an open problem: no physically-validated dataset existed even for this foundational task. The pipeline naturally extends to functional grasps by replacing the generation method with one targeting task-specific objectives (e.g., tool-use orientation), while the validation and selection stages remain unchanged.

\noindent[\rthree] AutoDex is complementary to sim/DR, not a replacement. Large-scale sim data captures broad distributional priors; our physically-validated data provides high-fidelity labels for fine-tuning --- following the established paradigm of pretraining on large-scale data and fine-tuning on task-specific real data. This makes AutoDex empirically meaningful even at 1.6K trajectories.

\rssparagraph{[\rtwo] Dataset statistics}
51 objects in Table IV have 100 trials each (1,678 trajectories). The remaining 49 objects have 3 trials each as a held-out test set (detailed in supplementary). We will complete 100 trials for all 100 objects before camera-ready, scaling the dataset to $\sim$10K trajectories. Failure breakdown will be included in the revision.

{
  \small
  \bibliographystyle{plainnat}
  \bibliography{references}
}

\end{document}
